%%%%%%%%%%%%%%%%%%%%%%%%%%%%%%%%%%%%%%%%%%%%%%%%%%%%%%%&&&&
%                 File: SEDOPTICA.tex                     %
%                     VERSION: 1.0                        %
%                 Date: July 22, 2021                     %
%                                                         %
%           LaTeX template file for use with              %
%             Óptica Pura y Aplicada papers               %
%                                                         %
%   This file requires style file, opa_format.sty         %
%                                                         %
%       (c) 2021 - Sociedad Española de Óptica            %
%%%%%%%%%%%%%%%%%%%%%%%%%%%%%%%%%%%%%%%%%%%%%%%%%%%%%%%&&&&

% History
% Prepared by Fernando Quintero <fquintero@utb.edu.co> on 22/07/21.
% Modified by Andres Marrugo <agmarrugo@utb.edu.co> on 11/08/21.


%%%%%%%%%%%%%%%%%%%%%%%%%% PREAMBLE %%%%%%%%%%%%%%%%%%%%%%%%%%%%%%
\documentclass[10pt,letterpaper]{article}
\usepackage{opa_format}


\begin{document}
%%%%%%%%%%%%%%%%%%%%%%%% BEGIN DOCUMENT %%%%%%%%%%%%%%%%%%%%%%%%%% 
%% The \type{} command allows to enter one of the following document types:
% 1. Research Paper.
% 2. Introduction to Research.
% 3. Tutorial Paper.
% 4. Other.
\type{RESEARCH PAPER // INTRODUCTION TO RESEARCH// TUTORIAL PAPER// OTHER\\Section: DO NOT WRITE HERE}

%% The \titleng{} and \titlesp{} commands allows to enter the titles in English and Spanish respectively.
\titleng{Title (in English)}
\titlesp{Título (en español)}

%% The \author{} command allows to enter the names of the authors.
\author{A. Author$^{1,2,S^*}$, B. Author$^{2,S}$, C. Author, ...}
\affil{1. Affiliation\\
2. Affiliation}

%% The "authorinfo" environment allows to enter the emails and occupations with their respective commands.
\begin{authorinfo}
    \email{*}{correo@autor.com}
    \occupation{S}{miembro	de SEDOPTICA / SEDOPTICA member}
\end{authorinfo}

%% The \affil{} command allows to enter the affiliations of the authors.
\received{xx/xx/xxxx}
\accepted{xx/xx/xxxx}

%% The \doi{} command allows to enter the DOI of the document.
\doi{xxxxx}






%%%%%%%%%%%%%%%%%%%%%% ABSTRACT AND KEYWORDS %%%%%%%%%%%%%%%%%%%%%
\vskip -1.55pc
\begin{mdframed}[linewidth=1pt] 

%% Use \begin{abstract} <content> \end{abstract} for initialize the "abstract" environment.
\begin{abstract}
Write a brief abstract of the work in English in a single paragraph. THIS CHART SHOULD NOT EXCEED THIS PAGE
\end{abstract}

%% The \keywords{} command allows to enter the document keywords.
\keywords{}

%% Use \begin{resumen} <content> \end{resumen} for initialize the "resumen" environment.
\begin{resumen}
Only for Spanish speaking authors. 
\end{resumen}

%% The \pclaves{} command allows to enter the document keywords in spanish (palabras clave).
\pclaves{}
\end{mdframed}






%%%%%%%%%%%%%%%%%%%%%%%%% REFERENCES %%%%%%%%%%%%%%%%%%%%%%%%%%%%%
% Add references with BibTeX or manually.

%%%%%%%%%% If using BibTeX:
% \bibliography{sample}
% \bibliographystyle{opajnl}

%%%%%%%%%% If preparing manually:
\begin{thebibliography}{1}
\newcommand{\enquote}[1]{``#1''}

\bibitem{born1999principles}
M. Born, E. Wolf, Principles of optics: Electromagnetic theory of propagation, interference and diffraction of light. Cambridge, Cambridge University Press (1999).

\bibitem{adelson1985spatiotemporal}
E. H. Adelson, J. R. Bergen, ``Spatiotemporal energy models for the perception of motion,'' \href{https://doi.org/10.1364/JOSAA.2.000284}{J Opt Soc Am A} \textbf{2}(2), 284-299 (1985).

\bibitem{kech2015role}
M. Kech, P. Vrana, M. Wolf, “The role of topology in quantum tomography,”\url{http://arxiv.org/abs/1503.00506}

\end{thebibliography}




%%%%%%%%%%%%%%%%%%%%%%%%%%%%  BODY  %%%%%%%%%%%%%%%%%%%%%%%%%%%%%%
\section{Introduction}
The first section must be an introduction to the subject of the paper. The main body text uses Computer Modern font, 10 points size with single spacing. References should appear between brackets [1], and numerated in order of appearance in the text.

\bigskip

\noindent The main body of the paper can be written in English or in Spanish.

\bigskip

Footnotes (notes at the bottom of text pages) are not used. Footnotes should be incorporated into the text or set with references as an endnote.

\section{Section Title 2}
\subsection{Subsection Title 1: Equations}
Equations must be centred and numbered in order of appearance, as
% 
\begin{equation}
    -\frac{1}{a}+\frac{1}{a'}=\frac{1}{f'} \enspace.
    \label{eq:eq1}
\end{equation}
% 
All equations should be numbered in the order in which they appear and should be referenced  from within the main text as Eq.~\eqref{eq:eq1}, Eq. (2), and so on.

\paragraph{Inline math} Place all math in a proper math environment. For example, expression \mbox{$8\times 2 = 16$} should be set this way, \texttt{\$8$\backslash$times 2=16\$}, not this way, \texttt{8 \$$\backslash$times\$2=16}. Simple fractions for inline math should use parentheses when necessary to avoid ambiguity, for example, to distinguish between $1/(n-1)$ and $1/n-1$.  Exceptions to this are the proper fractions such as $\frac{1}{2}$, which are
better left in this form. Summations and integrals that appear within text such as $\frac{1}{2}{\sum }_{n=1}^{n=\infty} (n^2 - 2n)^{-1}$ should have limits placed to the right of the symbol to reduce white space.

% \newpage

\subsection{Subsection Title 2: Tables}
Tables must include	a caption. They	must be	numbered in	order of appearance and should use the standard \LaTeX environment. 	

%% For tables use \begin{table} <content> \end{table}.
\begin{table}[]
\centering
\caption{Caption}
\label{tab:tab1}
\begin{tabular}{@{}cc@{}}
\toprule
Polarizer    & Analyzer \\ \midrule
0$^{\circ}$  & 90       \\
45$^{\circ}$ & 45       \\
90           & 0        \\ \bottomrule
\end{tabular}
\end{table}

\subsection{Subsection Title 3: Figures}
Figures should be placed centred in the body of the manuscript. Standard \LaTeX{} environments should be used to place tables and figures:
\begin{verbatim}
\begin{figure}[htbp]
\centering\includegraphics[width=7cm]{osafig1}
\caption{Sample caption (Fig. 2, \cite{Yelin:03}).}
\end{figure}
\end{verbatim}

Figures in colour are accepted and free of charge. 

%% For figures use \begin{figure} <content> \end{figure}.
\begin{figure}[!ht]
\centering
\includegraphics[width=0.47\textwidth]{figures/fig01.png}
\caption{Caption.}
\label{fig:01}
\end{figure}



\subsection{Subsection Title 4: References}
References should be numbered consecutively in the order in which they are first referenced in the body of the paper. Set reference callouts with standard \verb+\cite{}+ command. 
To reference multiple articles at once, simply use a comma to separate the reference labels, e.g. \verb+\cite{born1999principles, adelson1985spatiotemporal, kech2015role}+, produces \cite{born1999principles, adelson1985spatiotemporal, kech2015role}.

When reference authors are mentioned in the text, use surnames only (unless further clarity is needed), and use "et al." and first author name when three or more authors are given.

\paragraph{Journal articles:} authors are comma-separated listed first (initial and surname), followed by the article's full title in quotes, the journal's title abbreviation, the volume number in bold, the issue number in Roman and parenthesis, inclusive page numbers, and the year in parentheses.

\paragraph{Book as a whole or book chapter:} authors or editors are comma-separated listed first (initial and surname), followed by title in italics, city and publisher and year in parenthesis. Page numbers may be added if applicable.

\paragraph{Internet links} should list the author, title and the full URL (universal resource locator). Include the date of access, if relevant.

\paragraph{Bib\TeX} may be used to create a file containing the references, whose contents (i.e., contents of \texttt{.bbl} file) can then be pasted into the bibliography section of the \texttt{.tex} file. A Bib\TeX{} style file, \texttt{opajnl.bst}, is provided.

If your manuscript already contains a manually formatted \verb|\begin{thebibliography}|... \verb|\end{thebibliography}| list, you should ensure that your manually-formatted reference list adheres to the OSA style accurately.

\section{Conclusions}
The last section must be the conclusions of the work.


\section*{Acknowledgments} 
Acknowledgments, if included, should appear at the end of the document. 
Use the command \verb+\section*{Acknowledgments}+  to create a nonnumbered section heading.

\end{document} 




%%%%%%%%%%%%%%%%%%%%%%%% ENG DOCUMENT %%%%%%%%%%%%%%%%%%%%%%%%%% 